\chapter{Einführung}
Das Wichtigste zuerst: Gib dem GitHub-Repository\footnote{\url{https://github.com/julianbht/bht-thesis-template-latex}} einen Stern!!! 

Aber im Ernst - die Idee ist, dass jeder ein bisschen zum Repository beiträgt. Wenn du eine Verbesserung hast, einen Fehler entdeckst oder auch nur einen Rechtschreibfehler findest, erstelle bitte einen Pull Request. Ziel ist es, dass alle gemeinsam beitragen, damit die Vorlage mit der Zeit immer besser wird.

Diese Vorlage richtet sich an Studierende der \gls{bht} und erzeugt am Ende ein \gls{pdf}-Dokument im offiziellen BHT-Layout.

Wenn du diese Anleitung zum ersten Mal liest, solltest du folgende Schritte abarbeiten:

\begin{enumerate}
    \item Erstelle einen Zeitplan mit einigen Meilensteinen und sende ihn Siamak über Discord.\footnote{\url{TODO: Discord Link}}

    \item Lies dir das FAQ durch.
\end{enumerate}

\section{Motivation}
Meine Motivation für diese Arbeit kommt zunächst einmal durch die schnelle und drastische Entwicklung von KI und dem Nutzen in der Lehre.
In meinem Praktikum bei Donista durfte ich generativ mit KI und dem Agenten Claude arbeiten.
Gerade in dieser Zeit kam aus meinem Umfeld viel Kritik an der KI-Nutzung in der Lehre auf, da viele Studenten KI für die Abgaben in den Übungen der jeweiligen akademischen Fächer nutzten und es größtenteils nicht gesichert war, einen Lerneffekt aus der Nutzung zu ziehen, sondern eher, mit wenig Arbeit durch den Kurs zu kommen.
Da wurde dann auch mein Interesse geweckt, da ich bis zum Praktikum im Unterricht kaum bis gar keine KI genutzt habe, im Praktikum dann aber sehr viel KI-Nutzung hatte, gerade im Bereich der Softwareentwicklung und des Harness Engineerings.
Das Interesse an Präsentationen ist dadurch entstanden, dass wir im Studium viele Präsentationen gemacht haben und ich mir überlegt hatte, wie ich meine Präsentationen verbessern kann. Im Studium habe ich mich mit meinen Kommilitonen getroffen und wir sind unsere Präsis einzeln durchgegangen, aber was ist mit Studierenden, die keine Lerngruppen haben oder für die es schwierig ist, das Ganze mit den ganzen verschiedenen Workflows der anderen zu managen?
Daher die Motivation, einen Präsentationscoach zu entwickeln, der Feedback zur Präsentationsart gibt. Der Hintergrund zum direkten und reflexiven Feedback ist daher gekommen, dass ich aktuell einen Wechsel in den Lehramtsweg plane und mir generell Gedanken gemacht habe, wie unterschiedlich die Lerneffekte von direktem und reflexivem KI-Feedback sind.

\section{Fragestellung}
Daher die Fragestellung, die sich daraus entwickelt hat:
\begin{center}
	\begin{minipage}{0.85\textwidth}
		\centering
		\textbf{Entwicklung und Evaluation eines KI-gestützten Präsentationscoaches: Vergleich von direktem und reflexionsorientiertem KI-Feedback bei studentischen Präsentationen}
	\end{minipage}
\end{center}

\section{Methode}

\section{Überblick}

